
\section{Literature review}

Recent literature on artificial intelligence for water infrastructure shows a progressive shift from hand-crafted control rules and simplified hydraulic models towards data-driven forecasting, optimisation and condition monitoring, particularly in systems dominated by pumping stations. Within this strand, contributions can be grouped, in first approximation, into:
\begin{itemize}
    \item[(i)] supervised learning for forecasting hydraulic states and demands;
    \item[(ii)] \acrfull{rl} for real-time control and scheduling;
    \item[(iii)] \acrfull{ml}-based predictive maintenance and energy optimisation;
    \item[(iv)] hybrid physics-guided surrogate modelling for instantaneous electrical-to-hydraulic prediction.
\end{itemize}
In addition to established forecasting, control and maintenance applications, recent work has begun to explore hybrid modelling strategies that combine simplified physical principles with data-driven components in order to emulate the instantaneous behaviour of pumping and hydraulic systems, often in the broader context of theory-guided machine learning and \acrfull{dt} technologies [Karpatne2017; Barkanyi2021; Ghorbani2025]. The following examples illustrate how these paradigms have been applied in contexts directly involving water pumping stations or closely related infrastructures.

\paragraph{Supervised learning for forecasting hydraulic states and demands}

A first line of work focuses on \acrshort{dl} models for forecasting hydraulic variables directly at pumping stations. Zhang et al.\ propose a data-driven framework for predicting water levels in cascade pumping stations of the South-to-North Water Diversion Project in China, where the rear pool of each station is hydraulically coupled with the forebay of the downstream station [Zhang2023]. Their model combines a one-dimensional \acrfull{cnn} with a \acrshort{lstm} network and augments this architecture with a self-attention mechanism and a bagging ensemble to better capture both local feature interactions and long-range temporal dependencies in the station’s operating data. Using historical measurements from nine pumping stations, the hybrid \acrshort{cnn}–\acrshort{lstm} with attention achieves higher accuracy than either standalone \acrshort{cnn} or standalone \acrshort{lstm}, with noticeable improvements in coefficient of determination and mean absolute error on water level prediction [Zhang2023]. In the context of “smart” pumping stations and digital twins, this work exemplifies how deep sequence models can replace or complement physically based calibration for key operational parameters, enabling early detection of unstable regimes and supporting alarm logic directly on predicted levels rather than only on raw measurements.\\ 

Closely related, \acrshort{dl} has been extensively explored for water demand forecasting in distribution systems, with a clear impact on pump operation planning. Zanfei et al.\ develop a multivariate \acrshort{lstm} model for short-term (hourly) urban water demand, explicitly designed to integrate both longer-term consumption history and short-term meteorological information [Zanfei2022]. Their architecture employs a dual-module structure: one \acrshort{lstm} branch processes a week-long sequence of past demand, while a second branch processes shorter sequences of weather variables; the latent representations are then merged to produce the final demand forecast [Zanfei2022]. Tested on a real water distribution system, the multivariate \acrshort{lstm} outperforms a conventional multilayer perceptron and a \acrfull{sarima} model, highlighting the benefit of combining climatic covariates with historical demand in a \acrshort{dl} framework [Zanfei2022]. The same study also reviews previous contributions where \acrshort{lstm} networks were used for hourly and daily urban water demand prediction, such as Mu et al.\ for Chinese district metering areas [Mu2020], consolidating the view that recurrent deep networks generally outperform classical linear time-series models in this domain. Within the context of pumping stations, these demand forecasting models are typically used as inputs to optimisation or control routines that determine which pumps should operate and at what speeds, given predicted demand and electricity prices.\\

\paragraph{\acrlong{rl} for real-time control and scheduling}

Beyond pure forecasting, \acrshort{rl} techniques have recently been introduced to learn pump operation policies directly from interaction with hydraulic simulators. Hu et al.\ formulate real-time pump scheduling in water distribution systems as a Markov decision process and apply a \acrfull{ppo} agent to control variable-speed pumps, with the objective of minimising energy costs while maintaining pressure and tank-level constraints [Hu2023]. Their reward function explicitly balances pump energy cost, hydraulic constraint violations and tank level variation; by carefully tuning the penalty structure for tank depletion, they report non-trivial energy savings relative to alternative penalty formulations while preserving adequate storage levels in the long term [Hu2023]. \\

A complementary example, closer to flood-control pumping, is given by Goswami et al., who study rainwater pumping stations in Tamil Nadu designed to protect flood-prone areas [Goswami2025]. They model the station and its retention basin in the Storm Water Management Model (SWMM) and train a deep \acrshort{rl} agent based on a \acrfull{ddqn} with a \acrfull{gru} as function approximator, using synthetic extreme rainfall events generated via the Huff method as training scenarios [Goswami2025]. The \acrshort{ddqn}–\acrshort{gru} controller simultaneously minimises the maximum basin water level and the number of pump switching operations; experiments show that the learned policy reduces maintenance-related switching and better controls basin levels compared with the existing rule-based strategy [Goswami2025]. Together, these studies illustrate how \acrshort{rl} can move from offline optimisation of pump schedules to online policies that adapt to stochastic inflows and dynamic electricity tariffs, a direction that is highly relevant for complex water-lifting systems with multiple pumps and storage elements.

\paragraph{\acrlong{ml}-based predictive maintenance and energy optimisation}

A further strand of research addresses trend prediction and alarm generation for pumping station operating parameters using \acrshort{ml} architectures that are not necessarily very deep. Shao et al.\ propose a trend prediction and operation alarm model for pumping station units based on \acrfull{mtl} combined with an attention mechanism [Shao2024]. Historical operating data (e.g.\ vibration, temperature, flow and pressure) are first compressed via \acrshort{pca} to extract common latent features that summarise the unit’s state; these features are then fed into a multi-task learning model that jointly predicts trends for several parameters, while an attention module dynamically re-weights the shared features according to their relevance under new operating conditions [Shao2024]. Using real data from a water pumping station, the proposed \acrshort{mtl} model with attention improves trend prediction accuracy and stability compared with single-task and static-weight baselines, and the authors use threshold analysis on the predicted statistics to design a multi-stage alarm system for the unit’s operating condition [Shao2024]. Although the underlying network is relatively shallow compared with modern deep architectures, this work is representative of “classical” \acrshort{ml} ideas—feature extraction, multi-task learning and attention—applied directly to pump-station supervisory data for predictive monitoring.\\

Condition monitoring and predictive maintenance for pumps constitute another area where artificial intelligence has been applied extensively, often with methods that are transferable to water-lifting installations. Hasan et al.\ present a fault diagnosis framework for centrifugal pumps that converts vibration time series into time–frequency scalogram images using continuous wavelet transforms and then classifies health states with an \acrfull{adcnn} [Hasan2021]. This approach leverages \acrshort{cnn}-based feature extraction on images, allowing automatic discrimination among multiple fault types and severity levels without manual feature engineering [Hasan2021].\\ 

Subsequent works build on this idea by merging multiple time–frequency representations and by experimenting with pre-trained \acrshort{cnn} backbones to further improve diagnostic accuracy on centrifugal pump test rigs. For example, Zaman et al.\ propose a dual-scalogram approach, where two complementary time–frequency images are processed through a convolutional autoencoder and an artificial neural network for robust fault classification, showing improved performance over single-representation baselines [Zaman2023; Zaman2024]. While many of these studies are conducted on laboratory pumps rather than full-scale water pumping stations, the same sensing modalities (vibration, acoustic, pressure) and similar fault modes (e.g.\ cavitation, bearing defects, impeller damage) make them directly relevant for industrial water-lifting applications.\\

\acrshort{ml}-based predictive maintenance has also been explored in adjacent sectors, with techniques that can be transferred to water infrastructure. Abdalla et al.\ analyse multivariate sensor streams from \acrfull{esp} systems in the oil and gas sector and design a predictive maintenance pipeline that combines \acrshort{pca} for dimensionality reduction with \acrfull{xgb} for failure classification [Abdalla2022]. Their framework is trained to detect failure events several days in advance, achieving classification performance (e.g.\ F1-scores around 0.7) that is sufficient to provide actionable early-warning capabilities for operators [Abdalla2022]. Even though the application domain differs from municipal water supply, the general pattern—multivariate sensor fusion, latent-state extraction via \acrshort{pca}, and gradient-boosted tree ensembles for classification—can be replicated for predictive maintenance of large pumps and motors in aqueduct or wastewater lifting stations, provided that similarly rich condition-monitoring data are available.\\

Finally, recent work explicitly targets energy optimisation in wastewater and water pumping systems using advanced \acrshort{dl}. Piri et al.\ develop a framework that combines a deep residual network (ResNet) with self-attention mechanisms and a Grey Wolf Optimizer (GWO) to model and minimise the energy consumption of variable-speed pumps in a wastewater pumping station serving the Zahedan wastewater treatment plant [Piri2025]. The ResNet–attention model learns complex non-linear relationships between operational variables (flow rates, heads, pump speeds) and energy use, while GWO is employed to optimise design parameters and operating strategies [Piri2025]. In simulation, the proposed approach achieves significant energy savings compared with traditional proportional–integral–derivative (PID) control and also outperforms alternative optimisation baselines, underlining the added value of very deep architectures with attention and meta-heuristic optimisation for pump energy management [Piri2025]. Although validated so far on a single wastewater facility, this framework exemplifies how \acrshort{dl} can be embedded in supervisory control and data acquisition (SCADA) environments to deliver near real-time optimisation of pump operation.\\



\paragraph{Hybrid physics-guided surrogate modelling for instantaneous electrical-to-hydraulic prediction}

A fourth and increasingly relevant class of approaches focuses on hybrid surrogate models that approximate the instantaneous transformation of electrical power into hydraulic output in pumping and hydraulic installations. In contrast to forecasting-oriented methods or control-oriented \acrshort{rl}, these models do not aim at predicting future states of the system. Instead, they seek to reproduce, at a given operating condition, the static or quasi-static functional relationship that links electrical input, pump configuration and hydraulic response. This perspective is closely related to theory-guided or physics-guided \acrshort{ml}, where partial physical knowledge is combined with flexible function approximators in order to improve generalisation and physical plausibility [Karpatne2017].

A first group of contributions in this area targets the performance and efficiency of centrifugal or lifting pumps using hybrid or grey-box models. Liu et al.\ propose a two-stage hybrid model for predicting the efficiency of a centrifugal pump at variable speeds, where the efficiency curve is segmented into low-flow and high-flow regions and different modelling strategies are used in each regime [Liu2022]. In the small-flow stage a data-driven soft sensor based on statistical learning is trained, while in the large-flow stage a simplified mechanistic representation is preferred, and this combination reduces sample requirements and improves extrapolation compared with purely empirical regressors [Liu2022].\\

Deng et al.\ develop a multistage hybrid model for performance prediction of centrifugal pumps that uses only rotational speed and valve opening as inputs and combines knowledge-based segmentation of operating regimes with data-driven components to obtain the full head and efficiency characteristics [Deng2022]. Huang et al.\ introduce a hybrid neural network where a theoretical hydraulic loss model is embedded into a feedforward network for energy performance prediction of centrifugal pumps, so that head, power and efficiency are jointly estimated from geometrical parameters and operating conditions [Huang2020].\\

Wang et al.\ compare different surrogate models, including response surface and Kriging, within an optimisation loop for centrifugal pump impeller design, and show that surrogate-assisted optimisation can increase hydraulic efficiency while reducing power consumption using only a limited number of computational fluid dynamics simulations [Wang2016]. Similar surrogate-assisted strategies are adopted in several turbomachinery optimisation studies, where radial basis function networks or genetic-algorithm-optimised backpropagation networks are trained as surrogates of high-fidelity hydraulic solvers and then used to explore the design space of multi-stage centrifugal pumps under multiple performance objectives [Zhao2023; Bellary2016]. Taken together, these works show that hybrid knowledge-and-data models can exploit limited operational or simulation data, such as speed, valve status and a small number of sampled operating points, to reconstruct pump performance and efficiency curves with higher robustness than purely black-box models.\\

A second line of work develops surrogate models that emulate computationally expensive hydraulic simulators, with the explicit goal of enabling real-time prediction or optimisation at the system scale. Mahmoodian et al.\ present a hybrid surrogate modelling strategy to simplify and accelerate a detailed urban drainage simulator, where a reduced-order representation is calibrated on input-output data generated by the original hydrodynamic model and then used as an emulator for scenario analysis [Mahmoodian2018].\\

Palmitessa et al.\ propose a physics-guided \acrshort{ml} surrogate that approximates a high-fidelity hydrodynamic model of urban drainage networks, trained on a limited set of rainfall-runoff simulations and capable of predicting flows and water levels with high coefficients of determination while reducing simulation times by one to two orders of magnitude [Palmitessa2022]. Woo et al.\ design a deep learning based surrogate for urban pluvial flooding, where a convolutional architecture is trained to reproduce the spatio-temporal inundation fields generated by a physics-based hydraulic model for a large set of synthetic storm events [Woo2025Flood].\\

Zhang et al.\ propose a \acrfull{gnn} based surrogate model that is trained to reproduce the output of a physics-based flow routing model for urban drainage networks and show that hydraulic states can be predicted in real time with accuracy comparable to the original solver but at a fraction of the computational cost [Zhang2024GNN]. Recent work by Zandsalimi et al.\ extends this paradigm to end-to-end \acrshort{gnn} surrogates that model stormwater systems directly from network topology and boundary conditions, again with the aim of supporting real-time flood forecasting and control [Zandsalimi2025]. \\ 

In the context of optimisation, Li et al.\ integrate a hydrodynamic model with an evolutionary multi-objective algorithm for multi-barrage flood control and pumping station operation, using surrogate management to reduce the number of full hydraulic simulations required [Li2024HydroOpt], while Faúndez-Lizama et al.\ introduce a feasibility-guided evolutionary algorithm where an \acrshort{ml} classifier acts as a feasibility predictor surrogate before expensive EPANET simulations are executed for pump station design and operation in a large water distribution network [Faundez2025]. These contributions illustrate how surrogate models can be coupled with hydraulic solvers in order to accelerate simulation and optimisation of pumping schemes while preserving the essential hydraulic constraints.\\

A third set of studies explicitly adopts the language of \acrfull{dt} for water distribution networks and related systems and often relies on data-driven surrogate models to replace or augment hydraulic equations. Mücke et al.\ develop a probabilistic \acrshort{dt} for leak localisation in water distribution networks, where a generative neural network is trained as a stochastic surrogate of the hydraulic model and used within a Bayesian inference framework to infer leak locations from pressure and flow measurements [Muecke2023].\\

Pandey et al.\ implement a \acrshort{ml}-based \acrshort{dt} for anomaly detection and localisation in a treated wastewater distribution network, combining linear and nonlinear classifiers to detect changes in pressure differences between sensor pairs and to estimate posterior probabilities of possible leak locations [Pandey2024]. Dui et al.\ propose a \acrshort{dt} architecture for smart water distribution networks that integrates hydraulic models, real-time sensing and artificial intelligence components, including surrogate models, for resilience evaluation and intelligent operation [Dui2025].\\

In parallel, Ghorbani Bam and co-authors review \acrshort{dt} applications in the water sector and emphasise the central role of surrogate and reduced-order models for achieving timely monitoring and optimised operation of complex water systems [Ghorbani2025]. Pedersen discusses the use of \acrshort{dt} concepts for urban drainage and highlights how fast surrogate components can support daily operation and emergency management in drainage systems with pumps and storage [Pedersen2025]. Sharan et al.\ design a \acrshort{dt} of a coastal aquifer system that couples a three-dimensional numerical groundwater model with surrogate ensemble models within a simulation–optimisation framework in order to identify optimal pumping patterns under saltwater intrusion constraints [Sharan2025]. These examples show that \acrshort{dt} concepts in the water sector are often underpinned by surrogate models that emulate slow hydraulic solvers or difficult-to-measure relationships. \\

Finally, there is a broader body of work on grey-box and hybrid digital twins in adjacent energy and thermal systems that, although not centred on water pumping, is conceptually aligned with hybrid physics-guided surrogate modelling. Gilani et al.\ present a grey-box and black-box combined model to optimise the energy performance of building systems, where simplified resistance–capacitance thermal networks are complemented with \acrshort{ml} components in order to capture unmodelled dynamics [Gilani2019]. Li et al.\ review and systematise grey-box models for building energy simulation and model predictive control and emphasise that such models achieve a favourable balance between physical interpretability and computational efficiency [Li2021GreyBox]. Evens et al.\ develop a \acrshort{dt} of a heat pump based on a neural-network grey-box model trained on data generated by a simplified thermodynamic model and show that this hybrid approach can emulate the behaviour of the system with high fidelity while remaining suitable for real-time use [Evens2025]. In a more generic setting, Ozdagli et al.\ introduce a physics-guided learning architecture where parameters extracted from physics-based simulations are injected into intermediate layers of a neural network surrogate in order to constrain the learning process and improve generalisation [Ozdagli2023]. These cross-domain examples reinforce the idea that hybrid surrogate models can bridge the gap between purely data-driven regressors and fully mechanistic simulators in complex energy conversion systems.\\

In the specific context of water-lifting plants, hybrid physics-guided surrogate models provide a natural framework for learning instantaneous mappings from electrical inputs and pump states to hydraulic outputs, such as discharge flows, heads and hydraulic power, while enforcing only weak physical consistency through approximate energy balance or modular pump representations. By representing each pump as a computational unit that receives electrical power and local state variables and by aggregating pump contributions in a neural model of the station manifold or collector, such surrogates can emulate the electrical-to-hydraulic transformation of an entire pumping station without explicitly solving the underlying hydraulic equations. This modelling paradigm is closely aligned with the objectives of the present work, which aims to construct a data-driven \acrshort{dt} of a water-lifting installation for instantaneous prediction of hydraulic performance from electrical and configuration data.
\paragraph{Conclusions of the literature review}

Taken together, the reviewed contributions show a coherent evolution of artificial-intelligence applications in industrial pump-dominated and hydraulically complex systems. Supervised \acrshort{dl} models for water level and demand forecasting have demonstrated clear benefits over classical time-series baselines and now play a central role in anticipatory pump operation planning [Zhang2023; Zanfei2022; Mu2020].\\

\acrshort{rl}-based approaches have moved pump scheduling from static optimisation to adaptive policies that can react to stochastic inflows and varying electricity tariffs while respecting hydraulic constraints [Hu2023; Goswami2025].\\ 

In parallel, \acrshort{ml}-driven trend prediction, fault diagnosis and energy optimisation have matured across clean-water, wastewater and related sectors, leveraging multivariate condition-monitoring data and vibration or acoustic measurements to support predictive maintenance and energy-aware operation [Shao2024; Hasan2021; Zaman2023; Zaman2024; Abdalla2022; Piri2025].\\

More recently, hybrid physics-guided surrogate models and \acrshort{dt} architectures have emerged as a complementary paradigm, using grey-box and theory-guided strategies to approximate hydraulic simulators, pump performance curves and network-scale behaviour in a computationally efficient way [Liu2022; Deng2022; Huang2020; Mahmoodian2018; Palmitessa2022; Zhang2024GNN; Muecke2023; Pandey2024; Faundez2025; Ghorbani2025].\\

While many of these studies remain confined to specific case studies, laboratory pump rigs or single facilities, the underlying methodologies, which include recurrent and convolutional neural networks for time-series and vibration data, attention-based multi-task models for multivariate \acrshort{scada} streams, deep \acrshort{rl} for pump scheduling and hybrid surrogate or \acrshort{dt} architectures for real-time hydraulic emulation, are increasingly transferable to large-scale water-lifting plants such as aqueduct pumping stations. The main prerequisites are the availability of high-quality operational data, adequate instrumentation for key hydraulic and electrical variables and at least coarse hydraulic information that can be used to guide or constrain the learning process. Within this landscape, hybrid physics-guided surrogate modelling appears particularly promising for applications that require fast, snapshot-level prediction of hydraulic performance from electrical inputs and pump configurations, which is precisely the setting addressed in the experimental part of this thesis.
